\chapterbackmatter{Solutions to Weinberg Exercises}
\ExerciseEditorialNote

\begin{enumerate}
  \WeinbergSolution{1}
  For the superpotential of Chapter 26, Problem 4,
  \[
    f=\Phi_1\Phi_3^2+\Phi_2(\Phi_3^2+a),
  \]
  the equations \(f_i=0\) have no common solution when \(a\ne0\).
  The vacuum energy is then positive and the zero-energy spectrum is
  empty, so
  \[
    \operatorname{Tr}(-1)^F=0\qquad(a\ne0).
  \]
  At \(a=0\), introduce
  \(X=(\Phi_1+\Phi_2)/\sqrt2\),
  \(Y=(\Phi_1-\Phi_2)/\sqrt2\), and \(Z=\Phi_3\).  Then
  \[
    f=\sqrt2\,XZ^2.
  \]
  The field \(Y\) is a free chiral multiplet.  Its zero-momentum bosonic
  and fermionic states pair, so its factor in the index vanishes; hence
  the full Witten index is again zero.  This does not imply broken
  supersymmetry.  Indeed \(Z=0\), with arbitrary \(X\) and \(Y\), solves
  every \(\mathcal F\)-flatness equation and gives zero energy.

  No gauge field is present, so there is no instanton mechanism that
  could generate a non-perturbative superpotential.  Perturbative
  non-renormalization preserves \(f=\sqrt2XZ^2\), while corrections to a
  nonsingular Kahler metric cannot give positive energy at a point where
  all \(f_i\) vanish.  Supersymmetry therefore cannot be broken by
  higher-order effects at \(a=0\), even though the ordinary Witten index
  is zero.

  \WeinbergSolution{2}
  With the vector of \(SO(N_c)\) normalized to have Dynkin index one,
  \[
    C_1=N_c-2,\qquad C_2=N_f,\qquad
    b_0=3(N_c-2)-N_f.
  \]
  Write the gauge-invariant symmetric meson superfield as
  \[
    M_{mn}=\sum_{a=1}^{N_c}\Phi_{am}\Phi_{an}.
  \]
  When the bare superpotential vanishes and \(N_f<N_c-2\), holomorphy,
  flavor symmetry, the anomalous chiral symmetry, and dimensional
  analysis allow only
  \[
    f_{\rm np}
      =\kappa
       \left(
         \frac{\Lambda^{\,3(N_c-2)-N_f}}{\det M}
       \right)^{1/(N_c-N_f-2)},
  \]
  where \(\kappa\) is numerical.  It has the required dimension three
  and produces a runaway rather than a vacuum at finite \(M\).
  For \(N_f=N_c-2\) or \(N_f>N_c-2\), the selection rules of
  Section 29.3 forbid a generated superpotential when none is present
  initially.  In every case the Wilsonian gauge kinetic function has
  its one-loop running term, while unrestricted nonholomorphic
  corrections may occur in \(D\)-terms.

  Now add an equal mass term
  \[
    f_{\rm bare}=\frac{m}{2}\operatorname{Tr}M.
  \]
  For \(N_f<N_c-2\), the full superpotential is
  \(f_{\rm bare}+f_{\rm np}\); its stationary points are isolated.
  After all vectors are integrated out, holomorphic scale matching gives
  \[
    \Lambda_{\rm low}^{\,3(N_c-2)}
      =m^{N_f}\Lambda^{\,3(N_c-2)-N_f}.
  \]
  Gaugino condensation in the remaining pure \(SO(N_c)\) theory
  therefore contributes
  \[
    f_{\rm low}
      =(N_c-2)\Lambda_{\rm low}^{3}
      =(N_c-2)
       \left[
         m^{N_f}\Lambda^{\,3(N_c-2)-N_f}
       \right]^{1/(N_c-2)},
  \]
  with the \(N_c-2\) choices of phase understood.  The same low-energy
  expression applies when \(N_f\geq N_c-2\): the mass spurion permits the
  holomorphic term that was forbidden at \(m=0\), and all matter is
  removed below \(\lvert m\rvert\).

  \WeinbergSolution{3}
  Varying the first-order action \eqref{eq:29.5.29} with respect to the
  unconstrained \(W_L\) gives Eq.~\eqref{eq:29.5.30}.  Equivalently,
  \[
    \widetilde W_L=-h'(\Phi)W_L.
  \]
  This superfield identity is the compact relation between all the
  component fields.  To display them explicitly, put
  \(\vartheta=\theta_L^{\mathsf T}\epsilon\theta_L\) and use
  \[
    W_L=\lambda_L+
      \frac12\gamma^\mu\gamma^\nu\theta_L f_{\mu\nu}
      -\ii\theta_LD+\vartheta\sl{\partial}\lambda_R
  \]
  together with
  \[
    h'(\Phi)
      =h'(a)-\sqrt2h''(a)
        (\theta_L^{\mathsf T}\epsilon\psi_L)
       +\vartheta\left[
          h''(a)\mathcal F
          -\frac12h'''(a)(\bar\psi_L\psi_L)
        \right].
  \]
  Matching equal powers of \(\theta_L\) first gives
  \[
    \widetilde\lambda_L=-h'(a)\lambda_L,
  \]
  and then
  \[
    \begin{aligned}
    \frac12\gamma^\mu\gamma^\nu\theta_L\widetilde f_{\mu\nu}
      -\ii\theta_L\widetilde D
    ={}&
      -h'(a)\left(
        \frac12\gamma^\mu\gamma^\nu\theta_L f_{\mu\nu}
        -\ii\theta_LD
      \right)
    \\
    &+\sqrt2h''(a)
      (\theta_L^{\mathsf T}\epsilon\psi_L)\lambda_L .
    \end{aligned}
  \]
  Finally, the quadratic component obeys
  \[
    \begin{aligned}
    \vartheta\sl{\partial}\widetilde\lambda_R
    ={}&
      -h'(a)\vartheta\sl{\partial}\lambda_R
      +\sqrt2h''(a)(\theta_L^{\mathsf T}\epsilon\psi_L)
       \left(
        \frac12\gamma^\mu\gamma^\nu\theta_L f_{\mu\nu}
        -\ii\theta_LD
       \right)
    \\
    &-\vartheta\left[
       h''(a)\mathcal F
       -\frac12h'''(a)(\bar\psi_L\psi_L)
      \right]\lambda_L .
    \end{aligned}
  \]
  In a purely bosonic background the bilinear corrections vanish, and
  the left-chiral part of the dual field strength is simply
  \((\widetilde f_{\mu\nu})_L=-h'(a)(f_{\mu\nu})_L\), with the
  right-chiral relation its complex conjugate.  Thus the real dual
  tensor mixes \(f_{\mu\nu}\) with its Hodge dual whenever \(h'(a)\) has
  both real and imaginary parts.
\end{enumerate}
